\chapter{Experimental Study}

This chapter presents the experimental configuration, the simulation environment,
the evaluation methodology, and the results obtained by the eleven evaluated SLAM
algorithms.
All experiments were conducted on a single dataset to ensure full comparability
between algorithms.

% ─────────────────────────────────────────────────────────────────────────────
\section{Experimental Setup}

\subsection{Robot Platform}

The experimental platform is the TurtleBot3 Burger, a small differential-drive
mobile robot designed for indoor navigation.
It has a maximum linear velocity of $0.22$\,m/s and carries its sensors and
onboard compute in a compact, lightweight frame.
The TurtleBot3 Burger is a widely used research platform and is representative
of the class of small ground robots that rely on SLAM for autonomous operation
in GPS-denied environments.

\subsection{Sensor Configuration}

The robot is equipped with four sensors that together cover the modalities
evaluated in this thesis.

\begin{itemize}
  \item 3D LiDAR. A rotating 3D LiDAR operating at $33.6$\,Hz
        (scan period $\approx 30$\,ms), producing dense point clouds.

  \item Stereo camera. A Leopard Imaging LI-AR0234CS-NVIDIA Hawk
        stereo camera with a $100$\,mm baseline, $1920 \times 1200$\,px
        resolution per lens, and a $121.5^{\circ}$ horizontal field of view, recording
        at $60$\,fps.
        A depth map at $640 \times 360$\,px is also provided.

  \item IMU. An inertial measurement unit integrated with the camera
        module, publishing linear acceleration and angular velocity at
        approximately $62$\,Hz.

  \item Wheel odometry. Pose increments derived from motor encoders.
\end{itemize}

\subsection{Software Framework}

All algorithms are executed under ROS\,2 Jazzy on the same hardware.
Sensor data are replayed from a recorded rosbag at real-time speed, ensuring
identical temporal conditions for every algorithm.

Two preprocessing steps were applied where required by specific algorithms.
For the 2D LiDAR SLAM methods (GMapping, Cartographer~2D, and SLAM~Toolbox),
the 3D point cloud was projected onto a planar 2D scan by retaining only points
within a fixed height band around the robot.
For LIO-SAM, each point in the cloud was assigned a ring index derived from its
elevation angle, as the algorithm requires an organised point cloud with
per-point ring information.
All other algorithms consumed the raw sensor topics without modification.

% ─────────────────────────────────────────────────────────────────────────────
\section{Simulation Environment}

\subsection{Environment Description}

\subsection{Dataset Recording}

\subsection{Ground Truth Extraction}

% ─────────────────────────────────────────────────────────────────────────────
\section{Evaluation Methodology}

\subsection{Trajectory Accuracy}

\subsection{Map Quality}

\subsection{Computational Cost}

% ─────────────────────────────────────────────────────────────────────────────
\section{Results}

% Auto-generated by scripts/collect_metrics.py
% Bag: bag\_6
\begin{table}[H]
  \centering
  \setlength{\tabcolsep}{4pt}
  \caption[Trajectory and map quality metrics N/A bag6]{Trajectory and map quality metrics N/A bag\_6. ATE and RPE in metres. Row shading: \colorbox{tier1}{\phantom{X}}~ATE$<$0.1\,m, \colorbox{tier2}{\phantom{X}}~0.1--2\,m, \colorbox{tier3}{\phantom{X}}~$>$2\,m.}
  \label{tab:metrics_bag6}
  \renewcommand{\arraystretch}{1.2}
  \resizebox{\linewidth}{!}{%
  \begin{tabular}{lSSSSSSSSS}
    \toprule
    \rowcolor{hdr}
    \textbf{Algorithm} & \textbf{ATE} & \textbf{ATE} & \textbf{ATE} & \textbf{RPE} & \textbf{RPE} & \textbf{CD} & \textbf{F@5} & \textbf{F@10} & \textbf{F@20} \\
    \rowcolor{hdr}
    & \textbf{RMSE} & \textbf{mean} & \textbf{max} & \textbf{RMSE} & \textbf{mean} & \textbf{(m)} & \textbf{cm} & \textbf{cm} & \textbf{cm} \\
    \midrule
    \cellcolor{tier1}GMapping & 0.015 & 0.013 & 0.055 & 0.002 & 0.000 & 2.258 & 0.005 & 0.041 & 0.113 \\
    \cellcolor{tier1}DLIO & 0.042 & 0.035 & 0.376 & 0.003 & 0.002 & 1.598 & 0.043 & 0.099 & 0.196 \\
    \cellcolor{tier2}Cartographer 2D & 0.280 & 0.153 & 1.059 & 0.023 & 0.003 & 0.103 & 0.627 & 0.794 & 0.867 \\
    \cellcolor{tier2}KISS-ICP & 0.505 & 0.229 & 3.098 & 0.023 & 0.017 & 0.570 & 0.053 & 0.232 & 0.449 \\
    \cellcolor{tier2}MOLA Odometry & 0.223 & 0.215 & 0.331 & 0.003 & 0.003 & 0.210 & 0.016 & 0.249 & 0.597 \\
    \cellcolor{tier1}RTABMap LiDAR & 0.004 & 0.004 & 0.014 & 0.002 & 0.002 & 0.049 & 0.598 & 0.951 & 0.999 \\
    \cellcolor{tier3}SLAM Toolbox & 4.942 & 4.447 & 7.430 & 0.013 & 0.003 & 2.990 & 0.003 & 0.025 & 0.133 \\
    \cellcolor{tier2}RTABMap RGB-D & 0.102 & 0.052 & 0.690 & 0.007 & 0.000 & {N/A} & {N/A} & {N/A} & {N/A} \\
    \cellcolor{tier3}FAST-LIO2 & 3.974 & 3.485 & 8.558 & 0.028 & 0.007 & 2.756 & 0.003 & 0.046 & 0.133 \\
    \cellcolor{tier3}ORB-SLAM3 Stereo & 2.587 & 2.421 & 4.663 & 0.211 & 0.142 & {N/A} & {N/A} & {N/A} & {N/A} \\
    \cellcolor{tier3}ORB-SLAM3 S+IMU & 163.928 & 145.218 & 344.464 & {N/A} & {N/A} & {N/A} & {N/A} & {N/A} & {N/A} \\
    \cellcolor{tier2}cuVSLAM & 0.712 & 0.656 & 1.345 & 0.007 & 0.006 & 3.293 & 0.003 & 0.024 & 0.061 \\
    \cellcolor{tier1}lidarslam\_ros2 & 0.034 & 0.032 & 0.058 & 0.007 & 0.006 & 0.327 & 0.018 & 0.234 & 0.512 \\
    \cellcolor{tier1}EKF & 0.007 & 0.003 & 0.051 & 0.009 & 0.003 & 0.085 & 0.195 & 0.723 & 0.978 \\
    \bottomrule
  \end{tabular}}%
  \renewcommand{\arraystretch}{1}
\end{table}


\subsection{Trajectory Accuracy}

\subsection{Map Quality}

\subsection{Computational Cost}

% ─────────────────────────────────────────────────────────────────────────────
\section{Discussion}

\subsection{LiDAR-Based Methods}

\subsection{LiDAR-IMU Methods}

\subsection{Visual Methods}

\subsection{Failure Cases}
